\documentclass[10pt]{article}
\usepackage[T1]{fontenc}
\usepackage[utf8]{inputenc}
\usepackage[french]{babel}
\usepackage{amsmath,amssymb,mathrsfs,mathpazo}
\usepackage{geometry}
\geometry{a4paper,margin=2cm}
\usepackage{graphicx,array,tabularx,tabularray,colortbl,multirow,multicol,booktabs,enumitem,pifont,ulem,fancybox,diagbox,cancel}
\usepackage{pstricks,pstricks-add,pst-plot,pst-tree,pst-node,pst-text,pst-func,pst-math,pst-bspline,pst-3dplot,pst-eucl,pst-solides3d}
\let\pscalculate\undefined
\usepackage{pst-bezier}
\usepackage{tikz,pgfplots,tkz-tab}
\usetikzlibrary{arrows,arrows.meta,patterns,calc,positioning,trees,decorations.markings,shapes.geometric,shapes.misc,shapes.arrows,intersections,decorations.pathmorphing,decorations.pathreplacing,backgrounds,shadows,babel}
\pgfplotsset{compat=1.18}
\usepgfplotslibrary{fillbetween}
\usepackage[np]{numprint}
\usepackage{listings,makecell,pas-tableur,twoopt}
\let\barre\undefined
\usepackage[pstricks]{tablvar}
\definecolor{darkgreen}{RGB}{0,100,0}
\providecommand{\red}{\color{red}}
\providecommand{\blue}{\color{blue}}
\providecommand{\green}{\color{green}}
\providecommand{\black}{\color{black}}
\providecommand{\white}{\color{white}}
\providecommand{\gray}{\color{gray}}
\providecommand{\cyan}{\color{cyan}}
\providecommand{\magenta}{\color{magenta}}
\providecommand{\yellow}{\color{yellow}}
\providecommand{\violet}{\color{violet}}
\providecommand{\orange}{\color{orange}}
\providecommand{\euro}{\texteuro}
\providecommand{\EUR}{\texteuro}
\providecommand{\ieme}{\textsuperscript{e}}
\providecommand{\degres}{\ensuremath{^\circ}}
\providecommand{\degree}{\ensuremath{^\circ}}
\providecommand{\vv}[1]{\overrightarrow{#1}}
\providecommand{\Coord}[1]{#1}
\providecommand{\decosix}{}
\providecommand{\starredbullet}{\ensuremath{\bullet}}
\pagestyle{empty}
\setlength{\parindent}{0pt}
\usepackage{adjustbox}\DeclareUnicodeCharacter{2236}{:}\DeclareUnicodeCharacter{2190}{\ensuremath{\leftarrow}}
\providecommand{\R}{\mathbb{R}}
\providecommand{\N}{\mathbb{N}}
\providecommand{\Z}{\mathbb{Z}}
\providecommand{\Q}{\mathbb{Q}}
\providecommand{\D}{\mathbb{D}}
\providecommand{\C}{\mathbb{C}}
\providecommand{\e}{{\mathrm{e}}}
\providecommand{\vect}[1]{\overrightarrow{#1}}
\providecommand{\vectt}[1]{\overrightarrow{\mathrm{#1}}}
\providecommand{\vvt}[1]{\overrightarrow{\mathrm{#1}}}
\providecommand{\np}[2][]{{#2~#1}}
\providecommand{\vv}[1]{{\overrightarrow{#1}}}
\renewcommand{\d}{{\mathrm{d}}}
\providecommand{\degres}{{^\circ}}
\providecommand{\degre}{{^\circ}}
\providecommand{\blue}{\color{blue}}
\providecommand{\red}{\color{red}}
\providecommand{\enskip}{\quad}
\providecommand{\llbracket}{\lbrack\!\lbrack}
\providecommand{\rrbracket}{\rbrack\!\rbrack}
\providecommand{\euro}{{\text{€}}}
\providecommand{\no}{{\text{n°}}}
\providecommand{\interval}[2]{{[#1;#2]}}
\providecommand{\intervFO}[2]{{[#1;#2[}}
\providecommand{\dotfill}{\ldots}

\providecommand{\ds}{\displaystyle}
\providecommand{\textding}[1]{#1}
\providecommand{\Oij}{\(\mathrm{O};\vec{\imath},\vec{\jmath}\)}
\providecommand{\Oijk}{\(\mathrm{O};\vec{\imath},\vec{\jmath},\vec{k}\)}
\providecommand{\Ouv}{\(\mathrm{O};\vec{u},\vec{v}\)}
\providecommand{\intervFF}[2]{\ensuremath{[#1;#2]}}
\renewcommand{\euro}{\ifmmode\text{\texteuro}\else\texteuro\fi}
\DeclareUnicodeCharacter{00B0}{\ifmmode^{\circ}\else\textdegree\fi}
\DeclareUnicodeCharacter{20AC}{\ifmmode\text{\texteuro}\else\texteuro\fi}
% Énoncés : Éducation nationale. Transcription et corrigés : APMEP, Denis Vergès et contributeurs. Découpage et classement : CoopMaths, Sébastien Lozano, Éric Elter et contributeurs. Adaptation de la présentation : Maths Explorer.
% CC BY-SA : https://coopmaths.fr/www/about/
% https://coopmaths.fr/alea/static/bac/2022/tex/bac_2022_05_sujet1_madagascar_4.tex
\begin{document}
\noindent\textit{APMEP / CoopMaths. CC BY-SA. Présentation adaptée par Maths Explorer.}\par\medskip
\medskip

\begin{tabularx}{\linewidth}{|X|}\hline
\textbf{Principaux domaines abordés :} Géométrie dans l'espace\\ \hline
\end{tabularx}

\medskip

L'espace est muni d'un repère orthonormé \Oijk.

On considère les points A$(5~;~0~;~-1)$, B$(1~;~4~;~-1)$, C(1~;~0~;~3), D(5~;~4~;~3) et E(10~;~9~;~8).

\medskip

\begin{enumerate}
\item 
	\begin{enumerate}
		\item Soit R le milieu du segment [AB].
		
Calculer les coordonnées du point R ainsi que les coordonnées du vecteur $\vect{\text{AB}}$.
		\item Soit $\mathcal{P}_1$ le plan passant par le point R et dont $\vect{\text{AB}}$ est un vecteur normal. Démontrer qu'une équation cartésienne du plan $\mathcal{P}_1$ est :

\[x - y - 1 = 0.\]

		\item Démontrer que le point E appartient au plan $\mathcal{P}_1$ et que EA = EB .
	\end{enumerate}	
\item  On considère le plan $\mathcal{P}_2$ d'équation cartésienne $x - z - 2 = 0$. 
	\begin{enumerate}
		\item Justifier que les plans $\mathcal{P}_1$ et $\mathcal{P}_2$ sont sécants.
		\item On note $\Delta$ la droite d'intersection de $\mathcal{P}_1$ et $\mathcal{P}_2$.
		
Démontrer qu'une représentation paramétrique de la droite $\Delta$ est :

\[\left\{\begin{array}{l c r}
x&=&2+t\\
y&=&1+t \\
z&=&t\end{array}\right.(t \in \R).\]

	\end{enumerate}
\item  On considère le plan $\mathcal{P}_3$ d'équation cartésienne $y + z - 3 = 0$.

Justifier que la droite $\Delta$ est sécante au plan $\mathcal{P}_3$ en un point $\Omega$ dont on déterminera les coordonnées.
\end{enumerate}

\medskip

Si S et T sont deux points distincts de l'espace, on rappelle que l'ensemble des points M de l'espace tels que MS = MT est un plan, appelé plan médiateur du segment [ST]. On admet que les plans $\mathcal{P}_1$,\: $\mathcal{P}_2$ et $\mathcal{P}_3$ sont les plans médiateurs respectifs des
segments [AB], [AC] et [AD].

\begin{enumerate}[resume]
\item
	\begin{enumerate}
		\item Justifier que $\Omega \text{A} = \Omega \text{B} = \Omega \text{C} = \Omega \text{D}$.
		\item En déduire que les points A, B, C et D appartiennent à une même sphère dont
on précisera le centre et le rayon.
	\end{enumerate}
\end{enumerate}

\end{document}
